% ============================================================================
% Section 4 -- design. REWRITTEN 2026-09-07.
%
% The previous draft described POLICIES and never described the SYSTEM: a reader
% could not tell where a host-tier request's KV physically lives, what happens
% on a decode step, or how attention reads it. For a memory-system paper the
% architecture IS the contribution. It also narrated a superseded implementation
% ("our first version displaced residents"), which is lab notebook, not design.
%
% This version: architecture first, with a figure and a worked step; then the
% policies that drive it. Every mechanism here exists in the implementation --
% see the file/function names in section 5 -- so it can be checked against code.
% ============================================================================

\section{Design}
\label{sec:design}

\subsection{Overview}
\label{sec:design:overview}

The system presents device memory (HBM) and host memory (DRAM) as a single KV
store. Each request is assigned a \emph{tier}. A \textbf{device-tier} request is
served exactly as in stock vLLM: its KV blocks live in the paged device pool and
attention reads them there. A \textbf{host-tier} request holds \emph{no blocks
in the device pool at all}; its KV lives in host DRAM, and the bytes each layer
needs are streamed across the interconnect just before that layer's attention
consumes them.

Both tiers are served in the \emph{same} forward pass, so weights are read once
per step regardless of how the batch is split. Figure~\ref{fig:arch} shows the
four components.

\begin{figure*}[t]
\centering
\begin{tikzpicture}[
  x=1cm, y=1cm, font=\small,
  box/.style={draw, rounded corners=2pt, align=center, inner sep=2mm},
  mem/.style={box, fill=black!4},
  hot/.style={box, fill=blue!8, draw=blue!45},
  lbl/.style={font=\scriptsize\itshape, text=black!60, align=center},
  arr/.style={-{Latex[length=2mm]}, thick},
]

% ---------- host ----------
\node[mem, minimum width=3.2cm] (hk) at (0,0)
     {\texttt{host\_k}[$\ell$]\ \ \texttt{host\_v}[$\ell$]};
\node[lbl] (hkl) at (0,-0.85) {one pinned pair per layer,\\block-major, addressed by\\\textbf{host block id}};
\node[draw, dashed, rounded corners=3pt, fit=(hk)(hkl), inner sep=2.5mm] (hostbox) {};
\node[font=\bfseries\small] at (0,1.15) {Host DRAM};

% ---------- device ----------
\node[mem, minimum width=3.6cm] (pool) at (8.0,0.95)
     {paged KV pool\\[-1pt]{\scriptsize device-tier requests}};
\node[hot, minimum width=3.6cm] (roll) at (8.0,-0.95)
     {rolling region\\[-1pt]{\scriptsize $k$ slots $\times\ W$ blocks}\\[-1pt]{\scriptsize layer $\ell$ uses $\ell \bmod k$}};
\node[draw, dashed, rounded corners=3pt, fit=(pool)(roll), inner sep=2.5mm] (devbox) {};
\node[font=\bfseries\small] at (8.0,2.3) {GPU HBM};

% ---------- attention ----------
\node[box, minimum width=2.5cm] (attn) at (14.0,0)
     {attention\\layer $\ell$\\[-1pt]{\scriptsize one forward pass}};

% ---------- flows ----------
\draw[arr, blue!65] (hostbox.east|-hk) -- (4.4,0) -- (4.4,-0.95) -- (roll.west);
\node[lbl, text=blue!70] at (3.3,1.15)
     {\texttt{stream\_layer}($\ell$)\\refresh slot $\ell \bmod k$\\contiguous runs, H2D};

\draw[arr, black!60] (pool.east) -- (11.6,0.95) -- (attn.north west);
\node[lbl] at (11.55,1.55) {device\\block table};

\draw[arr, black!60] (roll.east) -- (11.6,-0.95) -- (attn.south west);
\node[lbl] at (11.55,-1.55) {compacted\\block table};

\draw[arr, red!60] (attn.south) -- (14.0,-2.9) -- (0,-2.9) -- (hostbox.south);
\node[lbl, text=red!70] at (7.0,-3.25)
     {\texttt{write\_token}($\ell$): this step's new $K,V$ $\rightarrow$ host mirror (authoritative) \emph{and} the current slot};
\end{tikzpicture}
\caption{Architecture. A host-tier request's KV lives in per-layer pinned host
mirrors addressed by an independent \emph{host block id} space. Once per step
the active cohort's blocks are compacted into a rolling region of $k$ device
slots; once per layer the slot for $\ell \bmod k$ is refreshed from the mirror.
Attention runs over both tiers in one forward pass, each tier reading its own
block table. The rolling region is sized to the \emph{active} cohort ($W$
blocks), not to the host pool, and is shared across all layers --- which is why
a host-tier request's device footprint is $O(1/k)$ of a resident one rather
than proportional to its context.}
\label{fig:arch}
\end{figure*}

\paragraph{Components.}
\begin{enumerate}
\item \textbf{Host KV pool.} Two pinned host tensors per layer, one for $K$ and
      one for $V$, laid out block-major. They are addressed by a \emph{host
      block id} drawn from a space of the system's own, unrelated to device
      block ids (\S\ref{sec:design:spaces}).
\item \textbf{Rolling region.} $k$ pairs of device tensors of $W$ blocks each.
      Layer $\ell$ reads slot $\ell \bmod k$, so $k$ slots cover the whole
      model and the device cost is independent of layer count.
\item \textbf{Tier-aware attention.} One forward pass; the batch's rows are
      partitioned by tier and each partition attends its own memory with its own
      block table.
\item \textbf{Tier manager.} Placement (\S\ref{sec:design:placement}), admission
      (\S\ref{sec:design:admission}) and migration
      (\S\ref{sec:design:migration}), plus the scheduler-side allocation rules
      that keep host-tier requests off the device pool.
\end{enumerate}

\subsection{Two address spaces}
\label{sec:design:spaces}

The single most consequential design decision is that the host mirror is
\textbf{not} indexed by device block id.

The reason is a lifetime mismatch. The point of the host tier is that a
host-tier request's device blocks can be released and reused by another
request. Under vLLM's allocator, releasing a block nulls the owner's block-table
entry and returns the id to a free list, from which it is promptly reissued. A
mirror indexed by device block id would therefore (i) lose the mapping needed to
find the request's own KV, since its block table no longer records where the
data was, and (ii) alias: the reissued id now denotes a different request's
data, so one request's stream would overwrite another's.

The system therefore assigns each \emph{(request, logical block)} a host block
id from its own space, held until the request finishes. Device ids may churn
freely underneath; the host mapping is unaffected. Growth needs no device
allocation either --- a new logical block is given a fresh host block, and its
tokens arrive through the per-layer write path rather than through a resident
write.

\subsection{The life of a decode step}
\label{sec:design:step}

Two granularities matter, and separating them is what keeps the per-layer work
small.

\paragraph{Once per step: compaction.} Given the active host-tier cohort, the
system takes the union of their host blocks, sorts it, and builds a bijection
onto $[0, W)$. This yields (i) the rolling region's size $W$ for this step, and
(ii) a \emph{compacted block table}: for each host-tier row, its logical blocks
expressed as rolling-region indices. The block table is built once and reused by
every layer, because the \emph{addresses} a request reads are fixed for the step
--- only the \emph{contents} of the slot change as layers advance.

Compaction also computes the maximal contiguous runs of host block ids in the
active set. A run is one contiguous byte range on both sides, so the per-layer
transfer is a handful of slice copies rather than a gather. This is why the host
mirror is block-major with $K$ and $V$ separated: with the block dimension
innermost, a run of blocks is strided, no pinned DMA is possible, and the
transfer degrades into a host-side copy.

\paragraph{Once per layer: refresh, write, attend.} For layer $\ell$:
\begin{enumerate}
\item \emph{Refresh.} Issue the runs from the host mirror into rolling slot
      $\ell \bmod k$. With $k > 1$ these can be issued a layer ahead, so the
      transfer overlaps the previous layer's compute.
\item \emph{Attend.} Device-tier rows attend the paged pool with the ordinary
      block table; host-tier rows attend slot $\ell \bmod k$ with the compacted
      block table. Both happen inside one forward pass over the shared weights.
\item \emph{Write.} The token generated this step projects new $K,V$. For
      host-tier rows these are written to the host mirror --- which is
      authoritative --- and into the current rolling slot, so that this step's
      own attention sees them.
\end{enumerate}

\paragraph{Why the device footprint is small.} A device-tier request's KV
occupies $L$ layers' worth of blocks. A host-tier request occupies host blocks,
plus its share of a rolling region that is $k$ slots deep regardless of $L$. Its
device footprint is therefore roughly $k/L$ of a resident request's: with
$k = 4$ and $L = 64$, about 6\%. That ratio, not the interconnect, is what
converts host DRAM into device capacity.

\subsection{Keeping host-tier requests off the device pool}
\label{sec:design:additive}

The capacity argument of \S\ref{sec:intro} requires host-tier requests to be
served \emph{in addition to} a full device-resident cohort rather than in place
of part of it. That holds only if a host-tier request never occupies a device
block, and three separate scheduler behaviours have to agree for that to be true.

\emph{Admission without allocation.} A host-tier request is admitted with an
empty block list. Its KV is delivered into the host pool by the prefill
transport (\S\ref{sec:background:pd}), and admission is contingent on the worker
confirming the KV is present --- the scheduler cannot observe a host-side write,
and admitting before it lands would leave a request with no KV anywhere.

\emph{No prefill allocation.} The request's prompt is marked already computed,
so no prefill is scheduled and no blocks are allocated even transiently. This
matters because the batch is bounded by \emph{peak} device residency: a design
that allocates during prefill and reclaims afterwards still pays the peak.

\emph{Queue bypass.} vLLM's waiting queue stops at the first request it cannot
allocate for. A host-tier request needs nothing, so a block-starved request
ahead of it would block an admission that requires no resources. Requests that
cannot be allocated are deferred so the scan continues.

For a request that is already resident when the system decides to move it, the
allocator additionally supports releasing its whole block list and, on the
return path, refilling it atomically. Release is triggered by a worker
acknowledgement that the KV reached the host, for the same reason admission is.

\subsection{Placement}
\label{sec:design:placement}

Which requests should be host-tier? The intuitive answer is the long ones, since
KV grows with context. \textbf{Context length cannot discriminate}, and the
reason is the cancellation of \S\ref{sec:model}: moving a request of context $S$
to the host frees $S\!\cdot\!k$ bytes of device memory and costs $S\!\cdot\!k$
bytes of interconnect on every subsequent step. Both scale identically, so
ranking by context is equivalent to ranking arbitrarily --- which
\S\ref{sec:eval:ablation} confirms.

The discriminating signal is \textbf{temporal}. A request that is not currently
generating holds device memory while consuming no interconnect bandwidth at all,
so the two costs come apart exactly when a request is idle or near the end of
its generation. Placement scores requests on estimated remaining generation and
time since last activity, and we keep the rule a transparent threshold so that
its behaviour can be explained; a learned variant is worth introducing only if
the simple rule leaves measurable headroom.

The rule uses only state the serving system already holds, so it needs no
externally supplied service class and applies in single-class deployments.

\subsection{Admission}
\label{sec:design:admission}

Each host-tier request costs $S\!\cdot\!k$ bytes of interconnect per step.
Admission enforces a budget on their sum: the step time multiplied by the
achieved interconnect bandwidth, less a margin.

Two properties of that budget matter more than its form. It is \textbf{measured,
not declared}: the step time comes from steps that carried no host-tier traffic.
A configured constant goes stale as the workload changes, and the running
latency average is circular, since host-tier traffic is exactly what it is being
used to bound. And the per-request cost is \textbf{derived from the model} ---
layer count, KV-head count, head dimension --- rather than configured, because a
configured value that disagrees with the running model is silently wrong in
whichever direction it disagrees.

A second, closed-loop guard sits behind the budget: an exponentially weighted
average of observed step time, compared against the latency objective with
hysteresis, gates new admissions when the system is already over. The budget
bounds what should fit; the guard responds to what did.

\subsection{Migration}
\label{sec:design:migration}

Load changes, so placement must be revisable: requests move to the host tier
under pressure and back when headroom appears, with hysteresis on both edges and
a cooldown bounding thrash.

The gate on those moves is measured rather than predicted, and this is a
deliberate departure. A cost model that asks whether a request's transfer fits
the remaining budget is \emph{necessary but not sufficient}, because the budget
accounts for transfer bytes while the host path also carries a per-step cost
independent of how many bytes move (\S\ref{sec:impl}). A gate that sees only
bytes cannot see that cost, and widening the margin does not repair the
omission.

A demotion is therefore a \textbf{probe}. The controller records the best recent
steady-state latency, performs the move, waits a settling period, and compares.
If latency regressed beyond a margin the request is promoted back and the
cooldown grows exponentially; otherwise the placement stands. Two details are
forced by the noise floor: the baseline must be a recent \emph{minimum} rather
than the instantaneous average, which would otherwise absorb startup transients,
and the loop needs a warmup so the first probe has a steady state to compare
against.

The probe is single-variable attribution in a system where requests arrive and
depart during the settling window, so its verdict is noisy; the margin and the
backoff bound the cost of a wrong one. Its value is that it is \emph{correct in
direction} where a byte-counting model is not, and \S\ref{sec:eval:ablation}
reports it reaching opposite verdicts on the two interconnects --- which is the
behaviour a placement policy should have when the hardware changes under it.
